\begin摘要}
针对高校地理科学专业GIS实践课"重操作轻思维"、与科研实践脱节的困境，本文设计并论述了基于生成式AI的城市物流园区选址分析教学案例，构建"问题定义—数据获取—空间建模—成果表达"四阶段人机协同教学模式。该模式以AI为技术协作者、教师为教学设计者，形成"学生—AI—教师"互动闭环，重点培养学生AI辅助GIS空间分析能力及科研创新能力。教学评价采用过程性评价（40\%）、成果性评价（40\%）与反思性评价（20\%）相结合的综合模式。创新点包括：从"操作导向"向"思维导向"的范式转变、"双师协同"教学模式重构、人机协同素养培育。本研究为培养适应新时代需求的复合型地理专业人才提供了可借鉴的教学模式。
\end摘要}

\begin关键词}
生成式AI；GIS实践课；人机协同；地理科学专业；教学创新